\section{泛型}

\begin{question}
	以下关于Java中的泛型的说法错误的是
	\begin{options}
		\item 可以提高代码的安全性
		\item 可以减少类型转换
		\item 泛型类型在运行时会被擦除
		\item 泛型可以用于基本数据类型
	\end{options}
	\begin{answer}
		D
	\end{answer}
	\begin{explanation}
		泛型不支持基本类型，只能用包装类。
	\end{explanation}
\end{question}

\begin{question}
	以下哪个是正确的Java泛型声明？
	\begin{options}
		\item class MyClass<T extends Number>
		\item class MyClass<T super Number>
		\item class MyClass<T implements Number>
		\item class MyClass<T>
	\end{options}
	\begin{answer}
		A
	\end{answer}
	\begin{explanation}
		泛型上限用 \texttt{extends}，不能用于类实现的接口（此处 Number 是类）。
	\end{explanation}
\end{question}
